\documentclass[11pt]{amsart}

% ---------------------------------------------------------------------------
%  The Lax--Phillips spectral criterion for the Riemann Hypothesis
%  A complete, self-contained problem specification.
%
%  Scope. This file specifies the CORRECT RH-equivalent Lax--Phillips
%  criterion (spectral-bound form) on the Sz.-Nagy--Foias model space
%  K_Theta attached to Theta(s) = xi(2s)/xi(-2s). It is the object that
%  REPLACES the unconditionally false energy-decay form (Lax eq. (54));
%  see Remark 1. Everything is defined from Hardy-space primitives and the
%  completed zeta function; two classical operator-theoretic facts are
%  isolated as Lemmas 1--2 with justification/citation, and the arithmetic
%  equivalence (Theorem 1) is proved in full.
%
%  Compile: pdflatex lax_phillips_rh_criterion.tex
% ---------------------------------------------------------------------------

\usepackage[utf8]{inputenc}
\usepackage[T1]{fontenc}
\usepackage{amsmath, amssymb, amsthm, mathtools}
\usepackage[margin=1in]{geometry}
\usepackage{microtype}
\usepackage{enumitem}
\usepackage{xcolor}
\usepackage[colorlinks=true, linkcolor=blue!50!black, citecolor=blue!50!black, urlcolor=blue!50!black]{hyperref}

\theoremstyle{plain}
\newtheorem{theorem}{Theorem}
\newtheorem{proposition}{Proposition}
\newtheorem{lemma}{Lemma}
\newtheorem{corollary}{Corollary}
\theoremstyle{definition}
\newtheorem{definition}{Definition}
\theoremstyle{remark}
\newtheorem*{remark}{Remark}
\newtheorem*{remarks}{Remarks}

\DeclareMathOperator{\dom}{dom}
\DeclareMathOperator*{\esssup}{ess\,sup}
\newcommand{\C}{\mathbb{C}}
\newcommand{\R}{\mathbb{R}}
\newcommand{\N}{\mathbb{N}}
\newcommand{\Cp}{\mathbb{C}_{+}}
\newcommand{\Ht}{H^{2}(\Cp)}
\newcommand{\Hi}{H^{\infty}(\Cp)}
\newcommand{\Kt}{K_{\Theta}}
\newcommand{\ip}[2]{\langle #1, #2\rangle}
\newcommand{\norm}[1]{\lVert #1 \rVert}
\newcommand{\abs}[1]{\lvert #1 \rvert}
\renewcommand{\Re}{\operatorname{Re}}

\title{\bfseries The Lax--Phillips spectral criterion\\ for the Riemann Hypothesis\\[2pt]
\large A complete, self-contained problem specification}
\author{OpenAI Codex}
\date{20 July 2026}

\begin{document}
\maketitle

\begin{abstract}
Let $\xi$ be the completed Riemann zeta function and
$\Theta(s)=\xi(2s)/\xi(-2s)$ the associated inner function on the right half-plane
$\Cp=\{\Re s>0\}$. On the Sz.-Nagy--Foias model space $\Kt=\Ht\ominus\Theta\Ht$ the
compressed translation (Lax--Phillips) semigroup $Z(t)$ has a generator $B$ whose
spectrum is $\{-\bar\rho/2\}$ over the nontrivial zeros $\rho=\beta+i\gamma$ of $\zeta$.
We give a fully self-contained specification of the criterion
\[
  \mathrm{RH}\iff s(B)\le-\tfrac14 ,
\]
where $s(B)=\sup\{\Re\lambda:\lambda\in\sigma(B)\}$ is the spectral bound. Every symbol
is defined from Hardy-space and analytic-number-theoretic primitives; the classical
operator-theoretic facts are isolated as two lemmas, and the arithmetic equivalence is
proved in full. We emphasise (Remark~1) that this is a statement about the
\emph{location of the spectrum} of $B$, not about $\norm{Z(t)}$: the growth bound is
$\omega_0(B)=0$ unconditionally, so the energy-decay form $\omega_0(B)\le-\tfrac14$
(Lax's eq.~(54)) is false independently of RH.
\end{abstract}

\tableofcontents

% ===========================================================================
\section{Function-theoretic preliminaries}
% ===========================================================================

\begin{definition}[Half-plane]\label{def:halfplane}
$\Cp:=\{s\in\C:\Re s>0\}$ is the open right half-plane, with boundary
$i\R=\{iy:y\in\R\}$ and closure $\overline{\Cp}=\{\Re s\ge0\}$.
\end{definition}

\begin{definition}[Hardy space $H^{2}$]\label{def:hardy}
$\Ht$ is the space of holomorphic functions $f:\Cp\to\C$ with
\[
  \norm{f}^{2}:=\sup_{x>0}\frac{1}{2\pi}\int_{-\infty}^{\infty}\abs{f(x+iy)}^{2}\,dy<\infty .
\]
For such $f$ the boundary limit $f(iy):=\lim_{x\downarrow0}f(x+iy)$ exists for a.e.\ $y$
and in $L^{2}(\R,\tfrac{dy}{2\pi})$; the supremum above is attained as $x\downarrow0$, and
the inner product is
\[
  \ip{f}{g}:=\frac{1}{2\pi}\int_{-\infty}^{\infty}f(iy)\,\overline{g(iy)}\,dy .
\]
Equivalently, by the Paley--Wiener--Laplace theorem, $f\in\Ht$ iff
$f(s)=\int_{0}^{\infty}e^{-st}g(t)\,dt$ for some $g\in L^{2}(0,\infty)$, and then
$\norm{f}=\norm{g}_{L^{2}(0,\infty)}$. Thus $\Ht$ is a Hilbert space.
\end{definition}

\begin{definition}[Reproducing kernel]\label{def:kernel}
For $w\in\Cp$,
\[
  k_{w}(s):=\frac{1}{\,s+\bar w\,}\in\Ht,\qquad \ip{f}{k_{w}}=f(w)\quad(f\in\Ht).
\]
(Indeed $k_{w}$ is the Laplace transform of $t\mapsto e^{-\bar w t}$, and
$f(w)=\int_{0}^{\infty}e^{-wt}g(t)\,dt=\ip{g}{e^{-\bar w t}}_{L^{2}(0,\infty)}$.)
\end{definition}

\begin{definition}[Bounded holomorphic multipliers]\label{def:hinf}
$\Hi$ is the algebra of bounded holomorphic $\varphi:\Cp\to\C$, with
$\norm{\varphi}_{\infty}:=\sup_{\Cp}\abs{\varphi}=\esssup_{y}\abs{\varphi(iy)}$. Each
$\varphi\in\Hi$ acts by multiplication $M_{\varphi}f:=\varphi f$ on $\Ht$, a bounded
operator with $\norm{M_{\varphi}}=\norm{\varphi}_{\infty}$; its adjoint fixes kernels,
\begin{equation}\label{eq:adjkernel}
  M_{\varphi}^{*}k_{w}=\overline{\varphi(w)}\,k_{w}\qquad(w\in\Cp),
\end{equation}
since $\ip{f}{M_{\varphi}^{*}k_{w}}=\ip{\varphi f}{k_{w}}=\varphi(w)f(w)
=\ip{f}{\overline{\varphi(w)}k_{w}}$.
\end{definition}

\begin{definition}[Inner function]\label{def:inner}
$\Theta\in\Hi$ is \emph{inner} if $\abs{\Theta(iy)}=1$ for a.e.\ $y\in\R$ (equivalently
$\norm{\Theta}_{\infty}\le1$ with unimodular boundary values). Then
$\Theta\Ht:=\{\Theta h:h\in\Ht\}$ is a closed subspace of $\Ht$.
\end{definition}

\begin{definition}[Model space]\label{def:model}
For inner $\Theta$,
\[
  \Kt:=\Ht\ominus\Theta\Ht=\{f\in\Ht:\ip{f}{\Theta h}=0\ \ \forall h\in\Ht\},
\]
a closed subspace with orthogonal projection $P_{\Kt}:\Ht\to\Kt$. If $\Theta(a)=0$ for
some $a\in\Cp$, then $k_{a}\in\Kt$, because
$\ip{\Theta h}{k_{a}}=\Theta(a)\,h(a)=0$ for all $h\in\Ht$.
\end{definition}

\begin{definition}[$C_{0}$-semigroup, generator, spectral and growth bounds]\label{def:semigroup}
A family $\{Z(t)\}_{t\ge0}$ of bounded operators on a Hilbert space $\mathcal H$ is a
\emph{$C_{0}$-semigroup} if $Z(0)=I$, $Z(t+u)=Z(t)Z(u)$, and $t\mapsto Z(t)x$ is
continuous for every $x\in\mathcal H$. Its \emph{generator} is
\[
  Bx:=\lim_{t\downarrow0}\tfrac1t\bigl(Z(t)x-x\bigr),\qquad
  \dom B:=\{x:\text{the limit exists}\},
\]
and one writes $Z(t)=e^{tB}$. Define the \emph{spectral bound} and \emph{growth bound}
\[
  s(B):=\sup\{\Re\lambda:\lambda\in\sigma(B)\},\qquad
  \omega_{0}(B):=\lim_{t\to\infty}\tfrac1t\log\norm{Z(t)}=\inf_{t>0}\tfrac1t\log\norm{Z(t)}.
\]
For every $C_{0}$-semigroup one has $-\infty\le s(B)\le\omega_{0}(B)<\infty$; equality
holds when $B$ is normal, but may fail for non-normal $B$.
\end{definition}

% ===========================================================================
\section{The arithmetic input}
% ===========================================================================

\begin{definition}[Riemann zeta and completed zeta]\label{def:zeta}
For $\Re s>1$, $\zeta(s)=\sum_{n\ge1}n^{-s}$, continued meromorphically to $\C$ with a
single simple pole at $s=1$. The \emph{completed zeta function} is
\[
  \xi(s):=\tfrac12\,s(s-1)\,\pi^{-s/2}\,\Gamma\!\bigl(\tfrac s2\bigr)\,\zeta(s).
\]
\end{definition}

\noindent\textbf{Classical facts about $\xi$} (all unconditional):
\begin{enumerate}[label=(\roman*),leftmargin=2.2em]
  \item $\xi$ is \emph{entire of order $1$};
  \item \emph{functional equation:} $\xi(s)=\xi(1-s)$;
  \item \emph{reality:} $\xi(\bar s)=\overline{\xi(s)}$; in particular $\xi$ is real on
        $\R$, and $\xi(0)=\xi(1)=\tfrac12\neq0$;
  \item \emph{zeros:} the zeros of $\xi$ are exactly the nontrivial zeros of $\zeta$;
        every one lies in the open critical strip $0<\Re s<1$, and $\xi$ has no other
        zeros (the trivial zeros of $\zeta$ and the pole at $s=1$ are cancelled by the
        $\Gamma$- and $s(s-1)$-factors). Writing a zero as $\rho=\beta+i\gamma$ we have
        $0<\beta<1$, and the zero set is symmetric under $\rho\mapsto1-\rho$ (by (ii))
        and under $\rho\mapsto\bar\rho$ (by (iii)).
\end{enumerate}

\begin{definition}[Riemann Hypothesis]\label{def:rh}
$\mathrm{RH}$ is the assertion that every zero $\rho=\beta+i\gamma$ of $\xi$ has
$\beta=\Re\rho=\tfrac12$.
\end{definition}

% ===========================================================================
\section{The Lax--Phillips data}
% ===========================================================================

\begin{definition}[The inner function of $\xi$]\label{def:theta}
Define
\begin{equation}\label{eq:theta}
  \boxed{\ \Theta(s):=\frac{\xi(2s)}{\xi(-2s)}=\frac{\xi(2s)}{\xi(1+2s)}\ }
  \qquad\Bigl(\text{second form via }\xi(-2s)=\xi(1+2s)\Bigr).
\end{equation}
\end{definition}

\begin{lemma}[$\Theta$ is inner; its zeros]\label{lem:inner}
$\Theta$ is a meromorphic inner function on $\Cp$. Its zeros in $\Cp$ are exactly
\[
  \bigl\{\,a=\rho/2:\ \xi(\rho)=0\,\bigr\},\qquad \Re a=\beta/2\in(0,\tfrac12),
\]
each with the multiplicity of the corresponding $\rho$; on $\overline{\Cp}$ it has no
poles and no zeros other than these.
\end{lemma}

\begin{proof}
$\Theta$ is a ratio of entire functions, hence meromorphic on $\C$. Zeros arise from
$\xi(2s)=0\iff2s=\rho\iff s=\rho/2$ (real part $\beta/2>0$); poles from
$\xi(-2s)=0\iff s=-\rho/2$ (real part $-\beta/2<0$, in $\C\setminus\overline{\Cp}$). No
zero coincides with a pole, since $\rho/2=-\rho'/2$ would force $\rho=-\rho'$, impossible
for two zeros with real parts in $(0,1)$; and neither lies on $i\R$ (real parts
$\pm\beta/2\neq0$). Thus $\Theta$ is holomorphic and zero-free on $i\R$, and holomorphic
on $\overline{\Cp}$ apart from the zeros $\rho/2$.

\emph{Boundary modulus.} For $s=iy$ we have $-2s=\overline{2s}$, so by reality
(Definition~\ref{def:zeta}(iii)) $\xi(-2s)=\overline{\xi(2s)}$ and hence
$\abs{\Theta(iy)}=\abs{\xi(2iy)}/\bigl|\overline{\xi(2iy)}\bigr|=1$
(the value $2iy$ has real part $0$, so $\xi(2iy)\neq0$).

\emph{Interior bound.} $\Theta$ is of bounded type in $\Cp$ (a quotient of two functions
of bounded type). Its mean type is nonpositive: from
$\log\xi(w)=\tfrac w2\log\tfrac{w}{2\pi}-\tfrac w2+O(\log\abs{w})$ one obtains
\[
  \log\Theta(s)=\log\xi(2s)-\log\xi(1+2s)=-\tfrac12\log\tfrac{s}{\pi}+O\!\bigl(\tfrac{\log\abs{s}}{\abs{s}}\bigr),
  \qquad \Theta(s)\sim(\pi/s)^{1/2}\to0
\]
as $\Re s\to\infty$; in particular $\lim_{\sigma\to\infty}\sigma^{-1}\log\abs{\Theta(\sigma)}=0$.
A function of bounded type in $\Cp$ with mean type $\le0$ and boundary modulus $\le1$
satisfies $\abs{\Theta}\le1$ throughout (Smirnov/Nevanlinna maximum principle). Together
with $\abs{\Theta}=1$ on $i\R$, this makes $\Theta$ inner. (This is the standard inner
function attached to $\xi$; cf.\ \cite{nikolski,uetake}.)
\end{proof}

\begin{definition}[The multiplier semigroup]\label{def:multiplier}
For $t\ge0$ put
\[
  \varphi_{t}(s):=e^{-ts}\in\Hi,\qquad \norm{\varphi_{t}}_{\infty}=1,\qquad
  \varphi_{t+u}=\varphi_{t}\varphi_{u},\quad \varphi_{0}=1,
\]
using $\abs{e^{-ts}}=e^{-t\Re s}\le1$ on $\overline{\Cp}$. Under the Laplace
identification of Definition~\ref{def:hardy}, $M_{\varphi_{t}}$ is right translation by
$t$; it is the outgoing-translation form of the unitary evolution group in
Lax--Phillips scattering \cite{laxphillips,laxphillipsauto}.
\end{definition}

\begin{definition}[Model / Lax--Phillips scattering semigroup and its generator]\label{def:Z}
Define the compressed multipliers
\[
  A_{t}:=P_{\Kt}\,M_{\varphi_{t}}\big|_{\Kt}:\Kt\to\Kt\qquad(t\ge0),
\]
and the \emph{Lax--Phillips semigroup} as the model-side representative of Lax's
energy-space semigroup under the Fourier--Laplace unitary $\mathcal F$ (his convention
$\mathcal F\,Z(t)^{*}\mathcal F^{-1}=A_{t}$):
\begin{equation}\label{eq:Z}
  \boxed{\ Z(t):=A_{t}^{*}=P_{\Kt}\,M_{\varphi_{t}}^{*}\big|_{\Kt}\ }.
\end{equation}
Then $\{Z(t)\}_{t\ge0}$ is a $C_{0}$-semigroup of contractions on $\Kt$
($\norm{Z(t)}\le\norm{\varphi_{t}}_{\infty}=1$): the map
$\varphi\mapsto P_{\Kt}M_{\varphi}|_{\Kt}$ is the Sz.-Nagy--Foias $H^{\infty}$-functional
calculus of the model operator, an algebra homomorphism, so $A_{t}A_{u}=A_{t+u}$ and
hence $Z(t)Z(u)=Z(t+u)$; strong continuity is that of the model semigroup
\cite{sznagyfoias,nikolski}. Let $B$ be its generator, $Z(t)=e^{tB}$.
\end{definition}

\begin{lemma}[Spectrum of the generator]\label{lem:spectrum}
The reproducing kernels at the zeros of $\Theta$ are eigenvectors of $Z(t)$: for
$a=\rho/2$,
\begin{equation}\label{eq:eig}
  Z(t)\,k_{a}=A_{t}^{*}k_{a}=\overline{\varphi_{t}(a)}\,k_{a}=e^{-t\bar a}\,k_{a},
  \qquad\text{hence}\qquad B\,k_{a}=-\bar a\,k_{a}=-\tfrac{\bar\rho}{2}\,k_{a}.
\end{equation}
Moreover, by the Sz.-Nagy--Foias / Liv\v{s}ic spectral theorem for model operators,
\[
  \sigma(B)=\overline{\bigl\{-\tfrac{\bar\rho}{2}:\ \xi(\rho)=0\bigr\}}
  \ \cup\ \bigl\{\zeta\in i\R:\ \Theta\text{ has no analytic continuation across }\zeta\bigr\}.
\]
The second set is empty ($\Theta$ is meromorphic, hence continues analytically across all
of $i\R$), and the zeros $\rho/2$ do not accumulate at any finite point of $i\R$.
Consequently, with $\rho=\beta+i\gamma$,
\begin{equation}\label{eq:sB}
  \Re\!\bigl(-\tfrac{\bar\rho}{2}\bigr)=-\tfrac{\beta}{2},\qquad
  \boxed{\ s(B)=\sup_{\rho}\Bigl(-\tfrac{\beta}{2}\Bigr)=-\tfrac12\,\inf_{\rho}\beta\ }.
\end{equation}
\end{lemma}

\begin{proof}
Equation~\eqref{eq:eig} is \eqref{eq:adjkernel} together with $k_{a}\in\Kt$
(Definition~\ref{def:model}), so $\sigma_{p}(B)\supseteq\{-\bar\rho/2\}$. The spectral
characterisation of the model generator is classical
\cite[Part~B]{nikolski}, \cite{sznagyfoias}: the spectrum consists of the zeros of the
inner function inside $\Cp$ (translated by the generator, $a\mapsto-\bar a$) together
with its boundary singular support. Here $\Theta$ is a ratio of entire functions, so it
continues analytically across every point of $i\R$ (its poles lie in
$\C\setminus\overline{\Cp}$; Lemma~\ref{lem:inner}); and since the zeros $\rho/2$ have
$\abs{\gamma}\to\infty$ with only finitely many below any height, they have no finite
accumulation point on $i\R$. Hence $\sigma(B)=\overline{\{-\bar\rho/2\}}$, whose
supremum of real parts is $\sup_{\rho}(-\beta/2)=-\tfrac12\inf_{\rho}\beta$, giving
\eqref{eq:sB}.
\end{proof}

% ===========================================================================
\section{The criterion equivalent to RH}
% ===========================================================================

\begin{theorem}[Lax--Phillips spectral criterion for RH]\label{thm:crit}
Let $B$ be the generator of the Lax--Phillips semigroup $Z(t)$ on $\Kt$
(Definitions~\ref{def:theta}--\ref{def:Z}). Then
\[
  \boxed{\ \mathrm{RH}\quad\Longleftrightarrow\quad s(B)\le-\tfrac14
          \quad\Longleftrightarrow\quad \sigma(B)\subset\{\Re s\le-\tfrac14\}\ }.
\]
Moreover $s(B)\ge-\tfrac14$ unconditionally, so equivalently
$\mathrm{RH}\iff s(B)=-\tfrac14$: the spectrum of $B$ touches the line
$\Re s=-\tfrac14$ but never crosses it.
\end{theorem}

\begin{proof}
By Lemma~\ref{lem:spectrum}, $s(B)=-\tfrac12\inf_{\rho}\beta$. The set of real parts
$\{\beta\}$ is nonempty and symmetric about $\tfrac12$ (Definition~\ref{def:zeta}(iv)),
so $\inf_{\rho}\beta\le\tfrac12$; hence $s(B)\ge-\tfrac14$ unconditionally. Next,
\[
  s(B)\le-\tfrac14\iff \inf_{\rho}\beta\ge\tfrac12
  \iff \text{no zero has }\beta<\tfrac12 .
\]
By the $\rho\mapsto1-\rho$ symmetry, ``no zero with $\beta<\tfrac12$'' is equivalent to
``no zero with $\beta>\tfrac12$'', hence to ``every zero has $\beta=\tfrac12$'', i.e.\ to
$\mathrm{RH}$ (Definition~\ref{def:rh}). Combining with $s(B)\ge-\tfrac14$ gives the
final equivalence.
\end{proof}

\begin{remarks}
\begin{enumerate}[leftmargin=1.6em]
\item \textbf{It is a spectral, not a norm, statement.} The constrained quantity is the
spectral abscissa $s(B)$ (the location of the resonances $-\bar\rho/2$), \emph{not}
$\norm{Z(t)}$. Unconditionally $\norm{Z(t)}=1$ for every $t>0$: by
Sarason's theorem \cite{sarason}, $\norm{Z(t)}=\norm{A_{t}}
=\operatorname{dist}_{H^{\infty}}(\varphi_{t},\Theta\Hi)$, and the nontrivial zeros of
$\zeta$ cluster densely in imaginary part (Riemann--von Mangoldt density
$\sim\tfrac{1}{2\pi}\log T$; $\gg T\log T$ of them simple and on the critical line by
Conrey \cite{conrey}), which forces this distance to equal $1$ (a Montel--Hurwitz
argument on the clustered interpolation data). Hence the growth bound is
$\omega_{0}(B)=0$, and the energy-decay form
\[
  \lim_{t\to\infty}\tfrac1t\log\norm{Z(t)}\le-\tfrac14\qquad(\text{Lax's eq.~(54)})
\]
is \emph{false independently of RH}. Under RH one has the strict gap
$s(B)=-\tfrac14<0=\omega_{0}(B)$ (Definition~\ref{def:semigroup}), the signature of the
non-self-adjointness of $B$; identifying $s(B)$ with $\omega_{0}(B)$ is legitimate only
for normal generators and is exactly the step that fails here.

\item \textbf{Resonance reading.} In the half-plane variable the eigenvalue $-\bar\rho/2$
has real part $-\beta/2$, so the critical line $\Re\rho=\tfrac12$ maps to the resonance
line $\Re s=-\tfrac14$. The criterion says: all scattering resonances of the model lie on
$\Re s=-\tfrac14$.

\item \textbf{Convention-independence.} Replacing $Z(t)=A_{t}^{*}$ by $A_{t}$ replaces
$\{-\bar\rho/2\}$ by $\{-\rho/2\}$, with identical real parts $-\beta/2$; so $s(B)$ and
the criterion are unchanged, and $\norm{Z(t)}=\norm{A_{t}}$. The Fourier--Laplace sign is
fixed by contractivity ($\varphi_{t}=e^{-ts}$ decays on $\Cp$; $e^{+ts}$ is unbounded and
inadmissible).

\item \textbf{Scope of the identification.} The criterion is stated on the definite
Hilbert-space model $\Kt$. Whether Lax's energy space $E$ is unitarily---or merely
boundedly-invertibly (similarly)---equal to this $\Kt$ is the one external input to cite
(\cite{uetake}). Because $s(B)$ and $\omega_{0}(B)$ are \emph{similarity invariants}, the
criterion (and the refutation in Remark~1) transfer under any bounded-invertible
identification; a genuine unitary equivalence additionally gives the exact value
$\norm{Z(t)}=1$. The residual subtlety is definite-vs.-indefinite: the automorphic
Lax--Phillips energy form is a Krein/Pontryagin form, and Faddeev--Pavlov regularisation
modifies it; the specification above pins down the definite scalar realisation.
\end{enumerate}
\end{remarks}

% ===========================================================================
\begin{thebibliography}{9}

\bibitem{conrey}
J.~B.~Conrey,
\emph{More than two fifths of the zeros of the Riemann zeta function are on the critical line},
J.~Reine Angew.\ Math.\ \textbf{399} (1989), 1--26.

\bibitem{laxphillips}
P.~D.~Lax and R.~S.~Phillips,
\emph{Scattering Theory},
Pure and Applied Mathematics \textbf{26}, Academic Press, New York, 1967 (rev.\ ed.\ 1989).

\bibitem{laxphillipsauto}
P.~D.~Lax and R.~S.~Phillips,
\emph{Scattering Theory for Automorphic Functions},
Ann.\ of Math.\ Studies \textbf{87}, Princeton Univ.\ Press, 1976.

\bibitem{nikolski}
N.~K.~Nikolski,
\emph{Operators, Functions, and Systems: An Easy Reading}, Vols.\ 1--2,
Math.\ Surveys and Monographs \textbf{92--93}, Amer.\ Math.\ Soc., 2002.

\bibitem{sarason}
D.~Sarason,
\emph{Generalized interpolation in $H^{\infty}$},
Trans.\ Amer.\ Math.\ Soc.\ \textbf{127} (1967), 179--203.

\bibitem{sznagyfoias}
B.~Sz.-Nagy, C.~Foias, H.~Bercovici, and L.~K\'erchy,
\emph{Harmonic Analysis of Operators on Hilbert Space}, 2nd ed.,
Universitext, Springer, New York, 2010.

\bibitem{uetake}
Y.~Uetake, (2007).
Unitary identification of Lax's energy-space scattering semigroup with the compressed
translation (model) semigroup on $\Kt$.
\textup{[Reference as given in the audit brief; exact title/venue to be verified before
any submission.]}

\end{thebibliography}

\end{document}
